\section*{Capítulo \ref{ch:b3:spectral} --- Operadores compactos
e o teorema espectral}

\begin{solution}{exo:b3:spectral:1}
(a) $T(B)$ é um subconjunto limitado do espaço de dimensão finita
$\operatorname{im}T$: relativamente \omterm{def:b3:spectral:compact}{compacto} por Heine--Borel
(\cref{cor:b3:topology:heineborel}, transportado por um \omterm{def:b3:topology:continuity}{homeomorfismo}
linear com $\R^n$).
(b) $I$ \omterm{def:b3:spectral:compact}{compacto} significa que a bola unitária fechada é compacta, o que,
pelo teorema de Riesz (volume do segundo ano), acontece exatamente em dimensão
finita. Se um $T$ \omterm{def:b3:spectral:compact}{compacto} tivesse inversa limitada $T^{-1}$, então
$I = T^{-1}T$ seria \omterm{def:b3:spectral:compact}{compacto}
(\cref{prop:b3:spectral:ideal}): impossível em dimensão
infinita.
\end{solution}

\begin{solution}{exo:b3:spectral:2}
(a) $\norm{Tx}^2 = \sum\abs{\lambda_n}^2\abs{x_n}^2 \leq
\sup\abs{\lambda_n}^2\norm x^2$, com quase igualdade nos
$e_n$ que realizam o supremo.
(b) ($\Leftarrow$) As truncaturas $T_N$ (mantêm $n \leq N$,
zero além) têm posto finito e $\vertiii{T - T_N} =
\sup_{n>N}\abs{\lambda_n} \to 0$: \omterm{def:b3:spectral:compact}{compacto} pela~\cref{prop:b3:spectral:ideal}. ($\Rightarrow$) Se
$\abs{\lambda_{n_k}} \geq \delta > 0$ ao longo de uma subsequência:
$\norm{Te_{n_k} - Te_{n_l}}^2 = \abs{\lambda_{n_k}}^2 +
\abs{\lambda_{n_l}}^2 \geq 2\delta^2$: nenhuma subsequência
convergente de $(Te_{n_k})$.
(c) $T^* = $ diagonal com $(\bar\lambda_n)$: \omterm{def:b3:spectral:selfadjoint}{autoadjunto} se, e somente se,
todos os $\lambda_n \in \R$. Então a base canônica $(e_n)$ é uma
base ortonormal de autovetores, de autovalores $\lambda_n \to
0$: o teorema espectral literalmente.
\end{solution}

\begin{solution}{exo:b3:spectral:3}
Seja $(x_n)$ limitada. Então $(Bx_n)$ é limitada
($\vertiii B < \infty$); a \omterm{def:b3:topology:compact}{compacidade} de $S$ extrai
$SBx_{n_k} \to y$; a \omterm{def:b3:topology:continuity}{continuidade} de $A$ dá $ASBx_{n_k} \to
Ay$: $ASB$ é \omterm{def:b3:spectral:compact}{compacto}. Se $ST = TS = I$ com $T$ \omterm{def:b3:spectral:compact}{compacto} e
$\dim H = \infty$: $I = ST$ seria \omterm{def:b3:spectral:compact}{compacto}, contradizendo
o~\cref{exo:b3:spectral:1}(b) --- que é o mesmo enunciado
visto do outro lado.
\end{solution}

\begin{solution}{exo:b3:spectral:4}
(a) Por Cauchy--Schwarz em $y$: $\abs{T_kf(x)}^2 \leq
\bigl(\int\abs{k(x,y)}^2\dd y\bigr)\norm f_2^2$; integre em
$x$ (Tonelli): $\norm{T_kf}_2 \leq \norm k_{L^2(\square)}
\norm f_2$.

(b) A família $e_{mn}(x,y) = e_m(x)\overline{e_n(y)}$ é
ortonormal em $L^2(\intcc01^2)$ (Tonelli separa a integral
dupla). Total: se $h \perp$ a todos os $e_{mn}$, então, para cada
$m$, a função $y \mapsto \int h(x,y)\overline{e_m(x)}\dd
x$ (em $L^2$ por Cauchy--Schwarz e Tonelli) é ortogonal a
toda $\overline{e_n}$ --- e as conjugadas
$(\overline{e_n})$ formam uma \omterm{def:b3:hilbert:onb}{base hilbertiana} sempre que $(e_n)$ forma
(a conjugação é uma bijeção isométrica de $L^2$ que preserva
ortogonalidade e totalidade) ---, de modo que ela é $0$ q.t.p.; então, para q.t.p.\ $y$, $h(\cdot, y) \perp$ a todo
$e_m$: $h(\cdot, y) = 0$ q.t.p.: $h = 0$ (Tonelli). Logo
$(e_{mn})$ é uma \omterm{def:b3:hilbert:onb}{base hilbertiana}; expanda $k = \sum c_{mn}e_{mn}$.
A truncatura $k_N$ (índices $\leq N$) dá $T_{k_N}$ de
posto finito (imagem em $\operatorname{Vect}(e_1, \dots,
e_N)$) e, por (a),
\[
\vertiii{T_k - T_{k_N}} \leq \norm{k - k_N}_{L^2} \to 0 :
\]
$T_k$ é um limite em \omterm{def:b3:banach:operator}{norma de operadores} de posto finito: \omterm{def:b3:spectral:compact}{compacto}
(\cref{prop:b3:spectral:ideal}).
\end{solution}

\begin{solution}{exo:b3:spectral:5}
(a) $\lambda\norm x^2 = \langle x, Tx\rangle \geq 0$ num
autovetor. A fórmula é a~\cref{prop:b3:spectral:sanorm} com todos os valores $\langle x,
Tx\rangle \geq 0$: o valor absoluto é redundante.
(b) $(x, y) \mapsto \langle x, Ty\rangle$ é uma forma sesquilinear
hermitiana positiva (possivelmente degenerada); a demonstração usual de
Cauchy--Schwarz (expanda $\langle x + t\eu^{\iu\theta}y,
T(x + t\eu^{\iu\theta}y)\rangle \geq 0$ e tome o
discriminante) nunca usa a definitude.
\end{solution}

\begin{solution}{exo:b3:spectral:6}
(a) $\norm{Mf}_2 \leq \norm f_2$ e, em $f_n = \sqrt n\,
\mathbf 1_{\intcc{1 - 1/n}1}$ (vetores unitários), $\norm{Mf_n}_2
\geq 1 - \frac1n$: $\vertiii M = 1$; \omterm{def:b3:spectral:selfadjoint}{autoadjunto}, pois o
multiplicador é real. Autovalores: $xf(x) = \lambda f(x)$ q.t.p.\
força $f = 0$ q.t.p.\ fora do conjunto nulo $\{x = \lambda\}$: $f =
0$ em $L^2$.
(b) Com as mesmas $f_n$: $\norm{Mf_n - f_n}_2 \leq \frac1n \to
0$, ao passo que $f_n \rightharpoonup 0$ (para $g \in L^2$ fixo,
$\abs{\langle g, f_n\rangle} \leq \norm{g\,\mathbf
1_{\intcc{1-1/n}1}}_2 \to 0$ pelo TCD). Se $Mf_{n_k} \to h$ em
norma, então $f_{n_k} \to h$, forçando $h = 0$ (limite fraco), e ainda assim
$\norm h = 1$: nenhuma subsequência convergente.
(c) No~\cref{lem:b3:spectral:existence}, é precisamente a
extração ``$Tx_{n_k} \to y$'' que usa a \omterm{def:b3:topology:compact}{compacidade}; para $M$, as
sequências maximizantes se concentram perto de $x = 1$ e suas
imagens convergem fracamente para $0$, nunca em norma: o autovetor
no topo da imagem numérica simplesmente não existe.
\end{solution}

\begin{solution}{exo:b3:spectral:7}
(a) $V = T_k$ com $k(x, y) = \mathbf 1_{y < x} \in
L^2(\intcc01^2)$: \omterm{def:b3:spectral:compact}{compacto} (\cref{exo:b3:spectral:4}). Seu
adjunto é o operador de núcleo com núcleo $\overline{k(y,
x)} = \mathbf 1_{y > x}$: $V^*f(x) = \int_x^1f$; $V \neq V^*$
(teste em $f = \mathbf 1$).
(b) Se $Vf = \lambda f$, $\lambda \ne 0$: $Vf$ é \omterm{def:b3:topology:continuity}{contínua}
em $\intcc01$ (convergência dominada em $\int_0^x f$), de modo que $f
= \frac1\lambda Vf$ tem um representante \omterm{def:b3:topology:continuity}{contínuo}; então
$Vf$ é $\mathcal C^1$ (teorema fundamental do cálculo para
integrandos \omterm{def:b3:topology:continuity}{contínuos}), de modo que $f$ é $\mathcal C^1$, e
$\lambda f' = f$ com $f(0) = \frac1\lambda Vf(0) = 0$: $f =
C\eu^{x/\lambda}$ com $C = f(0) = 0$.
(c) $V$ é \omterm{def:b3:spectral:compact}{compacto} sem autovalor algum, exceto possivelmente
$0$ ($Vf = 0$ força $f = 0$ q.t.p., derivando a
integral --- de modo que nem mesmo $0$): a maquinaria espectral
realmente exige a \omterm{def:b3:spectral:selfadjoint}{autoadjunção}, não apenas a \omterm{def:b3:topology:compact}{compacidade}.
\end{solution}

\begin{solution}{exo:b3:spectral:8}
Escreva $x = \sum_ic_ie_i + z$, $z \in \ker T$, de modo que $\langle x,
Tx\rangle = \sum_i\mu_i\abs{c_i}^2$. \emph{Cota inferior}: na
esfera unitária de $V_k = \operatorname{Vect}(e_1, \dots,
e_k)$, $\langle x, Tx\rangle = \sum_{i\leq
k}\mu_i\abs{c_i}^2 \geq \mu_k$: o máximo em $V$ do mínimo
é $\geq \mu_k$. \emph{Cota superior}: sejam $\dim V = k$ e $W_k
= \overline{\operatorname{Vect}}(e_k, e_{k+1}, \dots) +
\ker T$, de codimensão $k - 1$ (seu \omterm{thm:b3:hilbert:decomposition}{complemento ortogonal} é
$V_{k-1}$); $V \cap W_k \neq \{0\}$ (uma aplicação linear $V \to
H/W_k \cong V_{k-1}$ de posto $\leq k - 1$ tem núcleo
não trivial), e um unitário $x \in V\cap W_k$ tem $\langle x,
Tx\rangle = \sum_{i \geq k}\mu_i\abs{c_i}^2 \leq \mu_k$: o
mínimo em $V$ é $\leq \mu_k$. Juntas: a primeira fórmula;
a segunda se demonstra simetricamente (teste $W = W_k$; para
$W$ arbitrário de codimensão $k-1$, $W \cap V_k \neq 0$ dá
um vetor unitário com $\langle x, Tx\rangle \geq \mu_k$).
Monotonicidade: $\langle x, Tx\rangle \leq \langle x,
Sx\rangle$ pontualmente se transfere por $\max\min$.
\end{solution}

\begin{solution}{exo:b3:spectral:9}
Reescreva $f - \lambda Tf = g$. Para $\lambda = 0$: $f = g$,
sempre unicamente \omterm{def:b3:groups:derived}{solúvel}. Para $\lambda \neq 0$: isso é
$(T - \frac1\lambda)f = -\frac g\lambda$ e, pela \omterm{thm:b3:spectral:fredholm}{alternativa de
Fredholm} (\cref{thm:b3:spectral:fredholm}) com os
autovalores $\lambda_n = \bigl((n + \frac12)\pi\bigr)^{-2}$ de
$T$ (\cref{ex:b3:spectral:minxy}): \omterm{def:b3:groups:derived}{solubilidade} única para todo
$g$ se, e somente se, $\frac1\lambda \neq \lambda_n$ para todo $n$, isto é,
\[
\lambda \neq \Bigl(n + \tfrac12\Bigr)^2\pi^2
\qquad (n = 0, 1, 2, \dots).
\]
Num $\lambda = (n+\frac12)^2\pi^2$ excepcional: existem soluções
se, e somente se, $g \perp \sin\bigl((n{+}\frac12)\pi x\bigr)$, e elas são então únicas a menos de somar múltiplos desse seno.
\end{solution}

\begin{solution}{exo:b3:spectral:10}
(a) $Sx = \lambda x$: comparando coordenadas, $0 = \lambda
x_1$ e $x_n = \lambda x_{n+1}$; se $\lambda \ne 0$, então
$x_1 = 0$ e, indutivamente, $x = 0$; se $\lambda = 0$, $Sx = 0$
força $x = 0$ ($S$ isométrico). Nenhum autovalor. $S^*x =
\lambda x$ se lê $x_{n+1} = \lambda x_n$: $x =
x_1(1, \lambda, \lambda^2, \dots)$, em $\ell^2$ exatamente quando
$\abs\lambda < 1$: um disco aberto inteiro de autovalores.
(b) $\norm{Se_n - Se_m} = \norm{e_{n+1} - e_{m+1}} = \sqrt2$:
a imagem da limitada $(e_n)$ não tem subsequência de Cauchy;
do mesmo modo, $S^*e_{n+1} = e_n$. Nenhum dos dois é \omterm{def:b3:spectral:compact}{compacto}.
(c) Para \omterm{def:b3:spectral:compact}{operadores compactos} \omterm{def:b3:spectral:selfadjoint}{autoadjuntos}, os autovalores
capturam tudo (\cref{thm:b3:spectral:spectral}); abandonando
qualquer das hipóteses, os autovalores podem sumir por completo ($S$,
Volterra) ou preencher um disco ($S^*$): o objeto robusto é o
espectro $\{\lambda : T - \lambda I \text{ não inversível}\}$,
cuja teoria pertence a um curso posterior.
\end{solution}

\begin{solution}{exo:b3:spectral:11}
(a) $S$ é o operador diagonal de coeficientes
$\sqrt{\mu_n} \to 0$: \omterm{def:b3:spectral:compact}{compacto}
(\cref{exo:b3:spectral:2}(b), transportado para a base
$(e_n)$ completada por $\ker T$, em que $S = 0$), \omterm{def:b3:spectral:selfadjoint}{autoadjunto}
(diagonal real), positivo ($\langle x, Sx\rangle =
\sum\sqrt{\mu_n}\abs{\langle e_n, x\rangle}^2$) e $S^2 =
T$ termo a termo.

(b) $R$ comuta com $T = R^2$; para um autovetor $x$ de
$T$ de autovalor $\mu$: $T(Rx) = RTx = \mu Rx$, de modo que o
autoespaço de dimensão finita $E_\mu = \ker(T - \mu)$ é
$R$-estável. Em $E_\mu$, $R$ é simétrico positivo com $R^2
= \mu\,\mathrm{id}$: seus autovalores $\rho$ satisfazem $\rho^2
= \mu$, $\rho \geq 0$: todos iguais a $\sqrt\mu$, e um
operador diagonalizável com um único autovalor é escalar:
$R = \sqrt\mu\,\mathrm{id}$ em $E_\mu$. Em $\ker T$: $\norm
{Rx}^2 = \langle x, R^2x\rangle = \langle x, Tx\rangle = 0$.
Logo $R$ coincide com $S$ em $\ker T$ e em todo autoespaço,
cujo espaço gerado fechado é $H$ (teorema espectral): $R = S$.

(c) $\sqrt G$ age como $\frac1{n\pi}$ em $e_n =
\sqrt2\sin(n\pi x)$: é o operador de núcleo com
\[
k(x, y) = \sum_{n\geq1}\frac{2\sin(n\pi x)\sin(n\pi
y)}{n\pi},
\]
convergindo a série em $L^2(\intcc01^2)$ (coeficientes
$\frac1{n\pi} \in \ell^2$; os núcleos das somas parciais dão as
aproximações de posto finito). Nenhuma forma fechada elementar é
necessária: o lado espectral \emph{é} o operador.
\end{solution}

\begin{solution}{exo:b3:spectral:12}
(a) $T^*T$ é \omterm{def:b3:spectral:compact}{compacto} (produto de um operador limitado por um
\omterm{def:b3:spectral:compact}{compacto}, \cref{exo:b3:spectral:3}), \omterm{def:b3:spectral:selfadjoint}{autoadjunto}
($(T^*T)^* = T^*T$), positivo ($\langle x, T^*Tx\rangle =
\norm{Tx}^2$). Núcleo: $T^*Tx = 0 \Rightarrow \norm{Tx}^2 =
\langle x, T^*Tx\rangle = 0 \Rightarrow Tx = 0$ e,
reciprocamente, $\ker T^*T = \ker T$. O teorema espectral
fornece os $(e_n)$ ortonormais com $T^*Te_n = s_n^2e_n$,
$s_n > 0$, que geram $(\ker T)^\perp$.

(b) $\langle f_m, f_n\rangle = \frac{\langle Te_m,
Te_n\rangle}{s_ms_n} = \frac{\langle e_m,
T^*Te_n\rangle}{s_ms_n} = \frac{s_n^2}{s_ms_n}\delta_{mn} =
\delta_{mn}$. Expanda $x = x_0 + \sum_n\langle e_n, x\rangle
e_n$ com $x_0 \in \ker T$ (\omterm{thm:b3:hilbert:parseval}{Parseval} no espaço gerado fechado mais o
núcleo); aplicando o \omterm{def:b3:topology:continuity}{contínuo} $T$:
\[
Tx = \sum_n\langle e_n, x\rangle\,Te_n =
\sum_ns_n\langle e_n, x\rangle\,f_n,
\]
convergindo a série porque suas somas parciais são de Cauchy
($\norm{\sum_{N<n\leq M}s_n\langle e_n, x\rangle f_n}^2 =
\sum s_n^2\abs{\langle e_n, x\rangle}^2$, dominado por
$\sup_{n>N}s_n^2\cdot\norm x^2$, e $s_n \to 0$).

(c) $\norm{Tx}^2 = \sum_ns_n^2\abs{\langle e_n, x\rangle}^2
\leq (\max s_n)^2\norm x^2$, atingido no $e_n$ maximizante: $\vertiii T = \max s_n$. Truncar a DVS no posto
$N$ deixa um operador de norma $\sup_{n>N}s_n \to 0$:
aproximação de posto finito. O operador de Volterra não tem
autovalores (\cref{exo:b3:spectral:7}), mas $V^*V$ tem:
$V^*Vf(x) = \int_x^1\int_0^tf(s)\,\dd s\,\dd t$ é um
operador de núcleo simétrico positivo (núcleo $1 - \max(x,y)$,
um núcleo de Green do tipo corda), cujos pares próprios --- calculados
via o problema de valores de contorno $-u'' = \lambda^{-1}u$, $u'(0)
= u(1) = 0$, isto é, a família do~\cref{exo:b3:spectral:9} --- dão
valores singulares $s_n = \bigl((n + \frac12)\pi\bigr)^{-1}$.
A DVS vive sobre \emph{duas} famílias ortonormais precisamente
porque $V$ gira sua geometria própria para longe: sem autovetores
e, ainda assim, com estrutura diagonal \omterm{prop:b3:galois:perfect}{perfeita} entre duas bases diferentes.
\end{solution}

\begin{solution}{pb:b3:spectral:1}
\textbf{1.} \omterm{def:b3:topology:continuity}{Continuidade}: $\min$ e $\max$ são \omterm{def:b3:topology:continuity}{contínuos};
simetria: trocar $x, y$ não troca nem $\min(x,y)$ nem $1 -
\max(x,y)$. Cotas: $0 \leq g$, e cotas do tipo $g(x,y) \leq
\max\cdot(1-\max)$ dão $g \leq \frac14$ (para $u
= \max$: $\min \leq u$, de modo que $g \leq u(1-u) \leq \frac14$). $g
\in L^2(\square)$: de Hilbert--Schmidt, logo $G$ \omterm{def:b3:spectral:compact}{compacto}
(\cref{exo:b3:spectral:4}); o núcleo é real simétrico: $G$
\omterm{def:b3:spectral:selfadjoint}{autoadjunto}.

\textbf{2.} Separando em $y = x$:
\[
u(x) = (1 - x)\int_0^x y\,f(y)\,\dd y +
x\int_x^1(1 - y)\,f(y)\,\dd y .
\]
Para $f$ \omterm{def:b3:topology:continuity}{contínua}, derive (produto e teorema
fundamental):
\[
u'(x) = -\int_0^xyf + \int_x^1(1-y)f
\qquad\text{(os termos de bordo se cancelam)},
\]
e $u''(x) = -xf(x) - (1 - x)f(x) = -f(x)$; claramente $u(0) =
u(1) = 0$. Reciprocamente, se $u \in \mathcal C^2$ se anula nos dois
extremos, $w = u - G(-u'')$ satisfaz $w'' = 0$, $w(0) = w(1) =
0$: $w$ é afim e se anula duas vezes, de modo que $w = 0$.

\textbf{3.} Seja $Gf = 0$, $f \in L^2$. Para $\psi \in \mathcal
C_c^\infty(\intoo01)$: $\psi = G(-\psi'')$ pela questão 2, de modo que
\[
\langle f, \psi\rangle = \langle f, G(-\psi'')\rangle
= \langle Gf, -\psi''\rangle = 0
\]
($G$ \omterm{def:b3:spectral:selfadjoint}{autoadjunto}): pelo lema fundamental
(\cref{cor:b3:lp:fundlemma}), $f = 0$ q.t.p. Positividade: para
$f$ \omterm{def:b3:topology:continuity}{contínua}, com $u = Gf$,
\[
\langle f, Gf\rangle = \int_0^1 fu = \int_0^1(-u'')u
= \bigl[-u'u\bigr]_0^1 + \int_0^1(u')^2 = \int_0^1(u')^2
\geq 0 ;
\]
para $f \in L^2$, aproxime em $L^2$ por $f_n$ \omterm{def:b3:topology:continuity}{contínuas}:
ambos os membros passam ao limite ($G$ limitado).

\textbf{4.} Se $Gu = \lambda u$, $\lambda \neq 0$: $Gu$ é
\omterm{def:b3:topology:continuity}{contínua} ($\abs{Gu(x) - Gu(x')} \leq \norm{g(x,\cdot) -
g(x',\cdot)}_2\norm u_2$, e o núcleo é uniformemente
\omterm{def:b3:topology:continuity}{contínuo}), de modo que $u$ tem um representante \omterm{def:b3:topology:continuity}{contínuo}; então
as fórmulas da questão 2 mostram que $Gu \in \mathcal C^2$, de modo que $u =
\frac1\lambda Gu \in \mathcal C^2$ com $-\lambda u'' =
-(Gu)'' = u$ e $u(0) = u(1) = 0$.

\textbf{5.} $-\lambda u'' = u$, $u(0) = 0$: $u = A\sin(x/\sqrt
\lambda)$ (positivo $\lambda$: pela questão 3, $\lambda =
\langle u, Gu\rangle/\norm u^2 > 0$ nos autovetores). $u(1) =
0$ força $\frac1{\sqrt\lambda} = n\pi$: $\lambda_n =
\frac1{n^2\pi^2}$, autofunções $\sin(n\pi x)$, normalizadas
$e_n = \sqrt2\sin(n\pi x)$ ($\int_0^12\sin^2(n\pi x)\dd x =
1$). Verificação de ortogonalidade: $2\sin(m\pi x)\sin(n\pi x) =
\cos((m-n)\pi x) - \cos((m+n)\pi x)$ integra a $0$ para $m
\neq n$.

\textbf{6.} $\ker G = \{0\}$ (questão 3), de modo que
o~\cref{thm:b3:spectral:spectral}(1) dá $H =
\overline{\operatorname{Vect}}(e_n)$: os senos são uma \omterm{def:b3:hilbert:onb}{base
hilbertiana} de $L^2(\intcc01)$. Para $f(x) = x(1 - x)$:
\[
c_n = \sqrt2\int_0^1x(1-x)\sin(n\pi x)\,\dd x
= \sqrt2\;\frac{2\bigl(1 - (-1)^n\bigr)}{n^3\pi^3}
= \begin{cases}\dfrac{4\sqrt2}{n^3\pi^3} & n \text{ ímpar},\\
0 & n \text{ par},\end{cases}
\]
(duas integrações por partes). \omterm{thm:b3:hilbert:parseval}{Parseval}: $\int_0^1x^2(1-x)^2\dd
x = \frac1{30} = \sum_{n \text{ ímpar}}\frac{32}{n^6\pi^6}$,
isto é, $\sum_{n\text{ ímpar}}n^{-6} = \frac{\pi^6}{960}$.

\textbf{7.} $\langle e_n, Ge_n\rangle = \lambda_n\norm{e_n}^2
= \lambda_n$. Para $x$ fixo, os coeficientes de $g(x,
\cdot)$: $\langle e_n, g(x,\cdot)\rangle = (Ge_n)(x) =
\lambda_ne_n(x)$, de modo que $g(x,\cdot) =
\sum_n\lambda_ne_n(x)\,e_n$ em $L^2$. A série explícita
$\sum_n\lambda_ne_n(x)e_n(y) = \sum_n\frac{2\sin(n\pi
x)\sin(n\pi y)}{n^2\pi^2}$ converge normalmente no quadrado
($\abs{\text{termo}} \leq \frac2{n^2\pi^2}$): sua soma é
\omterm{def:b3:topology:continuity}{contínua} e, para cada $x$, ela tem os mesmos
coeficientes em $L^2(\dd y)$ que $g(x, \cdot)$: as duas funções \omterm{def:b3:topology:continuity}{contínuas}
coincidem para todo $(x, y)$. Pondo $y = x$ e
integrando (a convergência normal permite a integração termo a
termo):
\[
\int_0^1g(x,x)\,\dd x =
\sum_n\lambda_n\int_0^1e_n(x)^2\dd x = \sum_n\lambda_n .
\]

\textbf{8.} $\int_0^1g(x,x)\dd x = \int_0^1x(1 - x)\dd x =
\frac16$, de modo que $\sum_{n\geq1}\frac1{n^2\pi^2} = \frac16$:
\[
\zeta(2) = \sum_{n\geq1}\frac1{n^2} = \frac{\pi^2}6 .
\]

\textbf{9.} Para $f = \mathbf 1$: $c_n =
\sqrt2\int_0^1\sin(n\pi x)\dd x = \sqrt2\,\frac{1 -
(-1)^n}{n\pi}$: $c_n = \frac{2\sqrt2}{n\pi}$ para $n$ ímpar, $0$
para par. \omterm{thm:b3:hilbert:parseval}{Parseval}: $1 = \sum_{n\text{
ímpar}}\frac{8}{n^2\pi^2}$, de modo que $\sum_{n\text{ ímpar}}n^{-2} =
\frac{\pi^2}8$, e $\zeta(2) = \frac{\pi^2}8\cdot\frac{1}{1
- \frac14} = \frac{\pi^2}6$ (os termos pares são $\frac14\zeta(2)$).
Comparação: \omterm{thm:b3:hilbert:parseval}{Parseval} para uma única $f$ soma
$\abs{\langle e_n, f\rangle}^2$; a fórmula do traço integra
a diagonal do \emph{núcleo}, o que equivale a somar
\omterm{thm:b3:hilbert:parseval}{Parseval} sobre uma família ortonormal inteira de uma só vez ---
$\sum_n\langle e_n, Ge_n\rangle$ --- e é, portanto, cega a
qualquer escolha particular de função de teste.

\textbf{10.} $\abs{c_n\cos(n\pi t)} \leq \abs{c_n}$ com
$\sum\abs{c_n}^2 < \infty$: para cada $t$ a série converge
em $L^2$ (expansão ortonormal,
\cref{thm:b3:hilbert:parseval}(3)); a cota das caudas
$\norm{u(t) - u_N(t)}_2^2 \leq \sum_{n>N}\abs{c_n}^2$ é
uniforme em $t$, e cada soma parcial é \omterm{def:b3:topology:continuity}{contínua} em $t$
(finitos cossenos): $t \mapsto u(t,\cdot)$ é \omterm{def:b3:topology:continuity}{contínua} com valores
em $L^2$. Para $f = \sum_{n\leq N}c_ne_n$: cada modo
$\cos(n\pi t)\sin(n\pi x)$ satisfaz $\partial_t^2 =
-n^2\pi^2 = \partial_x^2$ aplicado a ele, anula-se em $x = 0,
1$, tem valor $\sin(n\pi x)$ e derivada temporal $0$ em $t =
0$: a soma finita resolve tudo. Musicalmente: o
movimento da corda é uma superposição de \emph{ondas estacionárias} $e_n$,
cujas frequências $n\pi$ são a fundamental e seus
harmônicos; o teorema espectral diz que toda forma inicial
se decompõe unicamente nesses tons puros, sendo os coeficientes
$c_n$ o timbre. Ouvir uma corda é calcular uma
expansão ortonormal.

\textbf{11.} Com $a_n = \abs{\langle u_n, x\rangle}^2$:
$m_{p+1} = \sum_n\mu_n^{p+1}a_n = \sum_n\bigl(\mu_n^{p/2}
\sqrt{a_n}\bigr)\bigl(\mu_n^{p/2+1}\sqrt{a_n}\bigr) \leq
\sqrt{m_p\,m_{p+2}}$ (Cauchy--Schwarz em $\ell^2$). Logo
as razões $m_{p+1}/m_p$ são não decrescentes em $p$; como
$R(x) = \frac{m_1}{m_0}$, $\frac{\langle Ax, Ax\rangle}
{\langle x, Ax\rangle} = \frac{m_2}{m_1}$ e $R(Ax) =
\frac{m_3}{m_2}$, segue a cadeia, cada termo $\leq \mu_1$
porque $m_{p+1} \leq \mu_1m_p$ termo a termo. Convergência: se
$a_1 > 0$ (escrevendo o peso total do maior autovalor como
$a_1$), então
\[
\mu_1 \geq R(A^kx) = \frac{m_{2k+1}}{m_{2k}} =
\mu_1\,\frac{a_1 + \sum_{\mu_n<\mu_1}(\mu_n/\mu_1)^{2k+1}
a_n}{a_1 + \sum_{\mu_n<\mu_1}(\mu_n/\mu_1)^{2k}a_n}
\longrightarrow \mu_1,
\]
por convergência dominada das somas (razões $< 1$): o
método da potência converge para todo vetor inicial não
ortogonal ao autoespaço dominante.

\textbf{12.} Pontualmente, $\abs{\langle x, (A - B)x\rangle}
\leq \vertiii{A - B}\,\norm x^2$, de modo que $R_A(x) \leq R_B(x) +
\vertiii{A - B}$ para todo $x$. Alimentando com isso a
fórmula max--min do~\cref{exo:b3:spectral:8}: $\mu_n(A)
\leq \mu_n(B) + \vertiii{A - B}$ e, simetricamente em $A,
B$: $\abs{\mu_n(A) - \mu_n(B)} \leq \vertiii{A - B}$ para
todos os $n$ de uma vez.

\textbf{13.} $-w'' = x - x^2$ integra a $w = -\frac{x^3}6
+ \frac{x^4}{12} + cx$ (com $w(0) = 0$), e $w(1) = 0$
dá $c = \frac1{12}$:
\[
w = \frac{x^4 - 2x^3 + x}{12} = \frac{x(1-x)(1 + x -
x^2)}{12} = Gu .
\]
Então $\norm u_2^2 = \int_0^1x^2(1-x)^2 = \frac1{30}$ e,
com $\int_0^1x^3(1-x)^3 = B(4,4) = \frac1{140}$:
\[
\langle u, Gu\rangle = \frac1{12}\Bigl(\frac1{30} +
\frac1{140}\Bigr) = \frac{17}{5040},
\qquad
R(u) = \frac{17/5040}{1/30} = \frac{17}{168} .
\]
Logo $\frac1{\pi^2} = \lambda_1 \geq \frac{17}{168}$, isto é,
$\pi^2 \leq \frac{168}{17} = 9.8824$: $\pi \leq 3.14364 <
3.1437$ (valor verdadeiro $\pi^2 = 9.8696$). Um polinômio, uma
integral, um dígito.

\textbf{14.} Escreva $u = x - x^2$, de modo que $Gu =
\frac{u(1 + u)}{12}$ e, usando $\int u^2 = \frac1{30}$,
$\int u^3 = \frac1{140}$, $\int u^4 = B(5,5) =
\frac{4!\,4!}{9!} = \frac1{630}$:
\[
\norm{Gu}_2^2 = \frac1{144}\int u^2(1+u)^2 =
\frac1{144}\Bigl(\frac1{30} + \frac2{140} +
\frac1{630}\Bigr) = \frac1{144}\cdot\frac{62}{1260} =
\frac{31}{90720} .
\]
Logo $\frac{\langle Gu, Gu\rangle}{\langle u, Gu\rangle} =
\frac{31/90720}{17/5040} = \frac{31}{306}$ e, pela questão
11, isso ainda é $\leq \mu_1 = \frac1{\pi^2}$: $\pi^2 \leq
\frac{306}{31} = 9.87097$, isto é, $\pi \leq 3.14181 <
3.1419$ --- quatro dígitos (e o próximo iterado daria
cerca de oito, contraindo-se o erro por
$(\lambda_2/\lambda_1)^2 = \frac1{16}$ a cada passo).

\textbf{15.} A família $(e_m \otimes e_n)(x,y) =
e_m(x)e_n(y)$ é uma \omterm{def:b3:hilbert:onb}{base hilbertiana} de $L^2(\intcc01^2)$
(ortonormalidade por Tonelli; totalidade como no~\cref{exo:b3:spectral:5}). Pela questão 7, para $x$ fixo:
$g(x, \cdot) = \sum_n\lambda_ne_n(x)e_n$, de modo que o coeficiente
de $g$ em $e_m\otimes e_n$ é
\[
\langle e_m\otimes e_n,\ g\rangle
= \int_0^1 e_m(x)\,\lambda_n e_n(x)\,\dd x
= \lambda_n\,\delta_{mn} .
\]
\omterm{thm:b3:hilbert:parseval}{Parseval} no quadrado:
\[
\iint g^2 = \sum_{m,n}\abs{\langle e_m\otimes e_n,
g\rangle}^2 = \sum_n\lambda_n^2 .
\]

\textbf{16.} Pela simetria de $g$,
\[
\iint g^2 = 2\iint_{y<x}y^2(1-x)^2
= 2\int_0^1(1-x)^2\,\frac{x^3}3\,\dd x
= \frac23\,B(4, 3) = \frac23\cdot\frac{3!\,2!}{6!}
= \frac23\cdot\frac1{60} = \frac1{90} .
\]

\textbf{17.} Combinando: $\sum_n\frac1{n^4\pi^4} =
\frac1{90}$, isto é, $\zeta(4) = \frac{\pi^4}{90}$. Em
geral, $\operatorname{tr}(G^k) = \sum\lambda_n^k =
\frac{\zeta(2k)}{\pi^{2k}}$ é igual a uma integral iterada de
produtos do núcleo racional-polinomial $g$ sobre o
cubo de dimensão $k$: um número racional. Logo $\zeta(2k) \in
\pi^{2k}\Q$ para todo $k$. A máquina só alcança argumentos
pares porque os autovalores entram através de suas
\emph{potências} --- $\sum\lambda_n^k$ --- e $\lambda_n =
\frac1{n^2\pi^2}$: nenhuma combinação de traços produz $\sum
n^{-3}$; a natureza aritmética de $\zeta(3)$ (irracional por
Apéry, transcendência em aberto) fica além do inventário
espectral.

\textbf{18.} O maior termo de uma soma de termos positivos é no
máximo a soma: $\lambda_1^2 \leq \sum\lambda_n^2 =
\frac1{90}$, de modo que $\frac1{\pi^2} \leq \frac1{\sqrt{90}}$ e
$\pi \geq 90^{1/4} = 3.0801\ldots$ Com a questão 14:
$3.080 < \pi < 3.1419$, por pura aritmética de corda. Traços
mais altos afiam a cota inferior geometricamente:
$\lambda_1 \leq (\operatorname{tr}G^{2k})^{1/2k} =
\lambda_1\bigl(1 + \sum_{n\geq2}(\lambda_n/\lambda_1)^{2k}
\bigr)^{1/2k}$, e o fator parasita morre tão rápido quanto
$\bigl(\tfrac14\bigr)^{2k}\cdot\frac1{2k}$ --- o mesmo
mecanismo de lacuna espectral da convergência do método da potência
(questão 11), visto do lado dos traços.

\textbf{19.} $\abs{n^2\pi^2 - \nu} \geq \delta > 0$ para todo
$n$ (a sequência $n^2\pi^2 \to \infty$ evita $\nu$ com
folga), e $\abs{n^2\pi^2 - \nu} \geq \frac{n^2\pi^2}2$
para $n$ grande. Convergência $L^2$: os coeficientes
$\frac{c_n}{n^2\pi^2 - \nu}$ são somáveis ao quadrado (dominados
por $\frac{\abs{c_n}}\delta$). Convergência uniforme: as normas do sup
das caudas são limitadas por $\sqrt2\sum_{n>N}
\frac{\abs{c_n}}{\abs{n^2\pi^2 - \nu}} \leq
\frac{2\sqrt2}{\pi^2}\bigl(\sum\abs{c_n}^2\bigr)^{1/2}
\bigl(\sum_{n>N}n^{-4}\bigr)^{1/2} \to 0$ (Cauchy--Schwarz).
Verificação: $Gf + \nu Gu$ tem coeficiente em $e_n$
\[
\lambda_nc_n + \frac{\nu\lambda_nc_n}{n^2\pi^2 - \nu}
= \frac{c_n}{n^2\pi^2}\Bigl(1 + \frac{\nu}{n^2\pi^2 -
\nu}\Bigr) = \frac{c_n}{n^2\pi^2 - \nu} :
\]
exatamente os
coeficientes de $u$, de modo que $u = Gf + \nu Gu$; unicidade
porque uma diferença $v$ de soluções satisfaz $v = \nu
Gv$, isto é, $\langle e_n, v\rangle(n^2\pi^2 - \nu) = 0$ para
todo $n$: $v = 0$.

\textbf{20.} Como na questão 19, a equação $u = Gf + \nu
Gu$ é equivalente à família de equações nos coeficientes
$(n^2\pi^2 - \nu)\,\langle e_n, u\rangle = c_n$, $n \geq
1$. Para $n = m$ o membro esquerdo é
$0$: a \omterm{def:b3:groups:derived}{solubilidade} força $c_m = 0$, e então $\langle e_m,
u\rangle$ fica livre, ao passo que todos os demais coeficientes ficam
determinados: as soluções formam a reta $u_0 + \R e_m$.
Ressonância: uma forçante com componente no modo próprio
bombeia energia nele sem limite --- o balanço empurrado em
sua própria frequência.

\textbf{21.} Pela fórmula da questão 19, $\norm{R_\nu f}_2^2
= \sum\frac{\abs{c_n}^2}{(n^2\pi^2 - \nu)^2} \leq
\frac{\norm f^2}{(\pi^2 - \nu)^2}$ (para $\nu < \pi^2$ o
autovalor mais próximo é $\pi^2$), com igualdade aproximada em
$f = e_1$: norma do operador $\frac1{\pi^2 - \nu}$. \omterm{def:b3:topology:compact}{Compacidade}:
$R_\nu$ é o limite em norma de suas truncaturas de posto finito
(os coeficientes das caudas $\frac1{n^2\pi^2 - \nu} \to 0$);
a \omterm{def:b3:spectral:selfadjoint}{autoadjunção} e a positividade se leem na forma diagonal
(todos os coeficientes $\frac1{n^2\pi^2 - \nu} > 0$).
$R_\nu$ tem autovalores $\frac1{n^2\pi^2 - \nu}$: a
análise das Partes I--VI recomeça literalmente.

\textbf{22.} O dicionário: \emph{autovalor} $\lambda_n =
\frac1{n^2\pi^2}$ $\leftrightarrow$ frequência ao quadrado
$n^2\pi^2$ do $n$-ésimo harmônico; \emph{traço} $\sum
\lambda_n = \frac16$ $\leftrightarrow$ $\zeta(2) =
\frac{\pi^2}6$; \emph{norma de Hilbert--Schmidt} $\iint g^2 =
\frac1{90}$ $\leftrightarrow$ $\zeta(4) = \frac{\pi^4}{90}$;
\emph{\omterm{thm:b3:spectral:fredholm}{alternativa de Fredholm}} $\leftrightarrow$ ressonância da
corda forçada; \emph{min--max} $\leftrightarrow$
estimativas variacionais, até $\pi < 3.1437$ a partir de um único
polinômio. Por trás de cada emparelhamento, o mesmo objeto: um
\omterm{def:b3:spectral:compact}{operador compacto} \omterm{def:b3:spectral:selfadjoint}{autoadjunto}, diagonalizado uma vez e explorado
de cinco maneiras.

\textbf{23.} Separando por paridade e substituindo $n = 2m$
na parte par,
\[
\zeta(6) = \sum_{n\text{ ímpar}}\frac1{n^6} +
\sum_{m\geq1}\frac1{(2m)^6}
= \frac{\pi^6}{960} + \frac{\zeta(6)}{64},
\]
de modo que $\frac{63}{64}\zeta(6) = \frac{\pi^6}{960}$ e $\zeta(6)
= \frac{64\,\pi^6}{63\cdot960} = \frac{\pi^6}{945}$. O
\omterm{def:b3:topology:pathconnected}{caminho} do traço calcularia $\operatorname{tr}G^3 =
\sum\lambda_n^3 = \zeta(6)/\pi^6$ como $\iint g\,g_2$ com o
núcleo iterado $g_2(x,y) = \int_0^1g(x,z)g(z,y)\dd z$ ---
três integrações de polinômios definidos por partes; \omterm{thm:b3:hilbert:parseval}{Parseval} em
$x(1-x)$ precisou de apenas uma.

\textbf{24.} (a) Diagonalize: $v = \sum_na_nu_n$ (mais uma
possível componente no núcleo, com a qual $\langle v, Av\rangle$
nada ganha e $\norm v^2$ cresce, de modo que um maximizante não a
tem). Então $\langle v, Av\rangle = \sum\mu_na_n^2 \leq
\mu_1\sum a_n^2$, com igualdade se, e somente se, $a_n = 0$ sempre que
$\mu_n < \mu_1$: um maximizante está no
autoespaço de $\mu_1$. (b) Para qualquer $u$,
\[
\langle\abs u, A\abs u\rangle - \langle u, Au\rangle
= \iint k(x,y)\,\bigl(\abs{u(x)}\abs{u(y)} -
u(x)u(y)\bigr)\dd x\,\dd y \;\geq\; 0,
\]
sendo o integrando pontualmente não negativo. Se $P = \{u >
0\}$ e $N = \{u < 0\}$ têm ambos \omterm{def:b3:measure:measure}{medida} positiva, então, em
$P\times N$, o integrando vale $2k\abs{u(x)}\abs{u(y)} >
0$ num conjunto de \omterm{def:b3:measure:measure}{medida} positiva: desigualdade estrita. Uma
autofunção $u$ de $\mu_1$ maximiza o quociente de Rayleigh,
e $\abs u$ tem a mesma norma, de modo que a estritude exibiria
$R(\abs u) > \mu_1$ --- impossível; logo $u$ tem sinal
constante q.t.p., digamos $u \geq 0$. Então $u(x) = \mu_1^{-1}(Au)(x) =
\mu_1^{-1}\int k(x,y)u(y)\dd y > 0$ para todo $x \in
\intoo01$ ($k(x,\cdot) > 0$ e $u \neq 0$). (c) Se o
autoespaço tivesse dimensão $\geq 2$, ele conteria duas
autofunções \emph{ortogonais} $u, v$, cada uma de sinal
constante e sem zeros no \omterm{def:b3:topology:interior}{interior} por (b); mas então
$\abs{\langle u, v\rangle} = \int\abs u\,\abs v > 0$ ---
contradição. Na corda: $k = g > 0$ no quadrado aberto,
$\mu_1 = \lambda_1 = \frac1{\pi^2}$ é de fato
simples, $e_1 = \sqrt2\sin(\pi x)$ é positiva em
$\intoo01$; e cada $e_n = \sqrt2\sin(n\pi x)$, $n \geq 2$,
ortogonal à positiva $e_1$, deve integrar a zero
contra ela, e portanto muda de sinal --- como seus $n - 1$ zeros
\omterm{def:b3:topology:interior}{interiores} $\frac kn$ confirmam.

\textbf{25.} Como $n^2\pi^2 \to \infty$, o mínimo $d =
\min_n\abs{n^2\pi^2 - \nu}$ é atingido, em algum modo $m$,
e $d > 0$ porque $\nu$ evita o espectro. A fórmula
diagonal da questão 19 dá
\[
\norm{R_\nu f}_2^2 =
\sum_n\frac{\abs{c_n}^2}{(n^2\pi^2 - \nu)^2}
\leq \frac1{d^2}\,\norm f_2^2,
\]
com igualdade para $f = e_m$: $\vertiii{R_\nu} = \frac1d$,
o recíproco da distância de $\nu$ ao espectro
--- o princípio geral do resolvente, aqui em coordenadas
explícitas. A \omterm{def:b3:spectral:selfadjoint}{autoadjunção} se lê nos coeficientes diagonais
reais; a \omterm{def:b3:topology:compact}{compacidade} segue como na questão 21 (os
coeficientes tendem a $0$, de modo que as truncaturas de posto finito
convergem em norma). Preço da ressonância: para $f = e_1$ e $\nu =
(1-\varepsilon)\pi^2$, a fórmula dá $u =
\frac{c_1}{\pi^2 - \nu}e_1 = \frac1{\varepsilon\pi^2}e_1$,
contra a resposta estática $Ge_1 = \frac1{\pi^2}e_1$:
amplificação $\frac1\varepsilon$. Em $\varepsilon =
10^{-2}$ a resposta é $100$ vezes a estática --- e
ela diverge quando $\varepsilon \to 0$, o que é a alternativa da questão 20
vista do lado limitado: quanto mais próxima a frequência da
forçante de uma frequência natural, menos limitada a
inversa.
\end{solution}
